\section{Scope, Method, and Historical Coordinates}

\subsection{Question, corpus, and periods}

This study asks how diverse biblical, initiatory, pastoral, and ritual
practices became the modern Latin Church's restricted major exorcism, and what
the resulting history and current discipline do---and do not---authorize. Its
selected corpus begins with Israelite and Second Temple contexts, moves through
New Testament narratives, early Christian apology and initiation, late-antique
offices and representative medieval Western books, and then follows the Roman
Ritual, codification, twentieth-century reform, and current official sources.
It is not an exhaustive catalog of rites, demonologies, local practice, or
reported cases.

\subsection{Place, Churches, language, and terminology}

The historical account is principally Mediterranean and Western; its legal
center is the Latin Church, with U.S. territorial implementation considered
separately. A bounded comparison identifies why Latin law and ritual cannot be
projected onto Eastern Catholic Churches; it does not state Orthodox law or
practice. Ancient sources are read through the exact Greek, Latin, or identified
English witnesses recorded in the source audit, including the
Douay--Rheims/Challoner reception text, Robinson--Pierpont Byzantine Greek,
Charles's English \emph{Jubilees}, Whiston's Josephus, and identified patristic
and ritual editions. No Hebrew or Syriac wording supports a publication claim,
and absence of an original-language collation is stated where material.

``Exorcism'' is a family of historically changing uses, not the name of one
unchanged ceremony. ``Major exorcism'' here means the solemn Latin liturgical
act governed by canon 1172 and the applicable Roman ritual book. ``Minor
exorcism'' refers chiefly to authorized initiatory prayers and never by itself
implies possession. ``Deliverance'' creates neither a canonical category nor
permission. Doctrine, school theology, historical witness, universal law,
territorial implementation, clinical judgment, and project synthesis remain
distinct.

\subsection{Evidence classes, hierarchy, and method}

Canonical narratives are read in their own literary and theological settings.
Ancient and medieval texts establish what a witness says and how a practice was
represented, not an event's hidden cause or incidence. Official acts establish
their own law, rite, publication, or guidance; they do not retrospectively
authenticate a reported case. Modern scholarship assists reconstruction but
does not supply ecclesiastical authority. Clinical classifications and
safeguarding literature control claims within their professional scope, not
theological judgments.

The method separates contemporary or near-contemporary evidence, later
narrative and institutional memory, official reception, and modern
reconstruction. It records textual and jurisdictional variation rather than
silently harmonizing it. Named cases were excluded after a case-specific audit
did not recover a sufficient chain of contemporary, official, medical, and
reception evidence for the claims proposed.

\subsection{Legal scope and currentness}

The controlling universal law is the 1983 \emph{Codex Iuris Canonici} for the
Latin Church, promulgated by John Paul II, in its official Latin text,
especially canon 1172. Pertinent dicastery acts and the current Roman ritual
discipline are distinguished from the USCCB's account of the English edition
implemented for the Latin dioceses of the United States. The published
authentic-interpretations register and material amendments were checked through
26 July 2026; no change to canon 1172 was located. Particular law and diocesan
processes may add applicable controls and must be obtained from the competent
authority. This work gives no opinion on an individual case.

\subsection{Global uncertainties, rights, and review}

No auditable global dataset establishes the frequency or outcomes of major
exorcism. Ancient and medieval reports usually cannot be retrospectively
diagnosed, and several early liturgical corpora have disputed dates,
attributions, or compositional histories. The proprietary current ritual is
neither reproduced nor operationally summarized. Protected modern works are
quoted only in focused excerpts or paraphrased.

The historical-accounts profile governs the narrative. The articles profile
applies only to bounded theological-authority and current-law claims.
The expanded draft has received internal source, structural, and production
review, including every-page visual inspection. Independent theological,
canonical, patristic, historical-liturgical, clinical-safeguarding, pastoral,
rights, and production review records dated 27 July 2026 govern their stated
and deliberately bounded lanes. These reviews are not clinical validation,
professional medical advice, or ecclesiastical approval. Exact-snapshot
installation and distribution authorization remain separate gates. It is not
an official Church publication or an approved ritual, diagnostic, pastoral,
or legal guide.
